\documentclass[aspectratio=169]{beamer}

\usetheme{Madrid}
\usecolortheme{seahorse}

\usepackage{kotex}
\usepackage{graphicx}
\usepackage{amsmath}
\usepackage{booktabs}
\usepackage{array}
\usepackage{multirow}

\title[Learning a thousand tasks in a day]{Learning a thousand tasks in a day}
\author[K. Dreczkowski et al.]{
Kamil Dreczkowski \and Pietro Vitiello \and Vitalis Vosylius \and Edward Johns
}
\institute[Imperial College London]{
Robot Learning Lab\\
Imperial College London
}
\date[Sci. Robotics 2025]{Science Robotics, Vol. 10, eadv7594, 2025년 11월 12일}

\begin{document}

%------------------------------------------------
\begin{frame}
  \titlepage
\end{frame}

%------------------------------------------------
\begin{frame}{연구 개요}
  \small
  \begin{itemize}
    \item 분야: 로봇 조작 imitation learning, 다중 작업 학습, 데이터 효율성
    \item 핵심 문제: 작업당 시演 수가 매우 적을 때($<10$) 조작 정책 학습의 비효율성
    \item 제안 아이디어 1: 궤적을 정렬(alignment)과 상호작용(interaction) 두 단계로 분해하는 prior
    \item 제안 아이디어 2: 행동 복제(BC) 대신 시연 검색 기반(retrieval-based) 일반화 활용
    \item 제안 방법: MT3 (Multi-Task Trajectory Transfer) – 정렬/상호작용 모두 검색 기반으로 구성한 완전 분해형 imitation learning
    \item 주요 성과: 단일 시연만으로 1000개 일상 조작 작업 학습, 24시간 이내 사람 시연 시간으로 달성
  \end{itemize}
\end{frame}

%------------------------------------------------
\begin{frame}{연구 동기}
  \small
  \begin{itemize}
    \item 인간/동물: 수 회 시연만으로 조작 기술 습득 – 영아, 영장류, 설치류 사례
    \item 기존 로봇 BC 시스템: 작업당 수백–수천 시연 필요
    \begin{itemize}
      \item BC-Z: 100 작업에 약 26,000 시연
      \item RT-1: 744 작업, 약 130,000 시연
      \item MT-ACT: 38 작업, 7500 시연
    \end{itemize}
    \item 대규모 작업 수(수천 작업)로 확장 시, 시연 수집 비용과 인력 요구 과도
    \item 현실적 시나리오: 작업당 시연 수가 매우 제한된 few-shot / one-shot 조작 학습 필요
    \item 목표: 최소 시연으로도 많은 작업을 학습 가능한 구조적 prior와 검색 기반 일반화 탐구
  \end{itemize}
\end{frame}

%------------------------------------------------
\begin{frame}{핵심 공헌}
  \small
  \begin{itemize}
    \item 다중 작업 imitation learning을 few-demonstration-per-task 영역에서 체계적으로 평가
    \item 궤적 분해 + 검색 기반 일반화 결합한 MT3 제안
    \item 분해 기반 방법과 단일 정책(monolithic BC) 비교 – 데이터 크기와 작업 다양성에 따른 스케일링 분석
    \item MT3로 단일 시연만으로 1000개 일상 조작 작업 학습 및 402개 객체, 31 macro skill, 534 micro skill 평가
    \item MT3의 성공/실패 패턴, 개방루프(open-loop) 상호작용의 한계, 검색 및 포즈 추정 실패 모드 정량 분석
  \end{itemize}
\end{frame}

%------------------------------------------------
\begin{frame}{문제 정의: 작업, 스킬, 데이터 효율성}
  \small
  \begin{itemize}
    \item 목표: 단일 로봇이 다수 작업을 imitation learning으로 학습
    \item 작업 정의
    \begin{itemize}
      \item macro skill: open, insert, fold 등 상위 조작 primitive
      \item micro skill: 특정 객체 카테고리와 motion profile에 특화된 macro skill
      \item task: 특정 객체 인스턴스에서의 micro skill 실행
    \end{itemize}
    \item 시나리오: 작업당 시연 수가 매우 적은 few-shot / one-shot regime
    \item 비교 대상: 궤적 분해 기반 방법 4종 vs 단일 BC 정책(MT-ACT+)
    \item 평가 축: 시연 수(데이터 크기)와 작업 수/객체 수(데이터 다양성)에 따른 성공률과 스케일링 특성
  \end{itemize}
\end{frame}

%------------------------------------------------
\begin{frame}{시스템 개요 및 입력/출력 구조}
  \small
  \begin{itemize}
    \item 하드웨어: Sawyer 로봇 + Robotiq 2F-85 그리퍼
    \item 센서: 헤드 마운트 RealSense D415 RGB-D 카메라 1대
    \item 입력
    \begin{itemize}
      \item 타깃 객체의 분할된 포인트 클라우드
      \item 자연어 기반 작업 설명 문장
    \end{itemize}
    \item 출력: 다중 작업 정책이 생성하는 end-effector 제어 명령
    \item 공통 아키텍처: 모든 방법이 동일한 입력/출력 구조 사용, 차이는 정책 내부 설계(분해 여부, BC vs retrieval)
  \end{itemize}
\end{frame}

%------------------------------------------------
\begin{frame}{궤적 분해 개념}
  \small
  \begin{itemize}
    \item 기존 단일 정책: 정렬과 상호작용을 하나의 monolithic policy로 학습
    \item 제안 prior: 궤적을 두 단계로 분해
    \begin{itemize}
      \item 정렬 단계: end-effector 또는 grasped object를 타깃 객체에 대한 적절한 상대 pose로 이동
      \item 상호작용 단계: 실제 조작 궤적 실행 – 높은 정밀도 요구
    \end{itemize}
    \item 예시
    \begin{itemize}
      \item 정렬: 플러그를 소켓 앞에 위치시키는 동작 – 경로는 다양해도 최종 pose만 중요
      \item 상호작용: 플러그 삽입 궤적 – 미세한 궤적 차이도 성공에 영향
    \end{itemize}
    \item 기대 효과: 두 단계에 특화된 정책 학습으로 데이터 효율성 향상
  \end{itemize}
\end{frame}

%------------------------------------------------
\begin{frame}{정책 설계: 분해 기반 네 가지 조합}
  \small
  \begin{itemize}
    \item 각 단계별 두 가지 선택지
    \begin{itemize}
      \item 정렬: BC 기반 vs 검색 기반(pose estimation + motion planning)
      \item 상호작용: BC 기반 vs 검색 기반(open-loop trajectory replay)
    \end{itemize}
    \item 네 가지 조합
    \begin{itemize}
      \item BC-BC: 정렬과 상호작용 모두 BC
      \item BC-Ret: 정렬 BC + 상호작용 검색 기반
      \item Ret-BC: 정렬 검색 기반 + 상호작용 BC
      \item Ret-Ret(MT3): 정렬과 상호작용 모두 검색 기반
    \end{itemize}
    \item 비교 기준: 데이터 효율성, 일반화 성능, 스케일링 특성
  \end{itemize}
\end{frame}

%------------------------------------------------
\begin{frame}{정책 설계: Monolithic BC (MT-ACT+)}
  \small
  \begin{itemize}
    \item 단일 정책이 전체 궤적(정렬+상호작용)을 end-to-end로 학습
    \item 백본: MT-ACT 변형
    \begin{itemize}
      \item PointNet++ 기반 포인트 클라우드 인코더
      \item CLIP 임베딩 + FiLM으로 언어 조건부 특징 조정
      \item 변분 추론 기반 transformer – 다중 모달 행동 분포 모델링
    \end{itemize}
    \item 입력: 타깃 객체 포인트 클라우드(EE frame), 작업 설명, action history
    \item 출력: action chunk와 종료 신호
    \item 역할: 분해 기반 방법과 공정 비교 위한 강력한 monolithic BC 베이스라인
  \end{itemize}
\end{frame}

%------------------------------------------------
\begin{frame}{MT3: Multi-Task Trajectory Transfer 개념}
  \small
  \begin{itemize}
    \item Ret-Ret 조합에 해당하는 완전 검색 기반 분해 방법
    \item 핵심 구성
    \begin{itemize}
      \item 언어 + 기하 기반 시연 검색 – 테스트 전 전체 궤적 1개 선택
      \item 정렬: trajectory transfer 기반 pose estimation + motion planning
      \item 상호작용: 시연된 end-effector 속도 궤적을 EE frame에서 open-loop 재생
    \end{itemize}
    \item 특징
    \begin{itemize}
      \item 학습 시 네트워크 훈련 대신 시연 메모리와 검색 모듈 중심
      \item 테스트 시 시연을 그대로 재사용 – 행동이 항상 시연 궤적에 한정
    \end{itemize}
    \item 목표: 최소 시연으로 다수 작업 학습, 특히 one-shot 설정에서의 실용성 검증
  \end{itemize}
\end{frame}

%------------------------------------------------
\begin{frame}{시연 검색 파이프라인 개요}
  \small
  \begin{itemize}
    \item 2단계 계층적 검색 구조 (논문 p.12 그림 참조)
    \item 1단계: 언어 기반 검색
    \begin{itemize}
      \item 작업 설명에서 micro skill 이름 추출 (예: ``open bottle'')
      \item 동일 micro skill에 해당하는 시연 집합 필터링
    \end{itemize}
    \item 2단계: 기하 + pose 기반 검색
    \begin{itemize}
      \item PointNet++ 인코더로 포인트 클라우드 임베딩
      \item 코사인 유사도로 test 객체와 가장 유사한 시연 선택
      \item 전역 형상 유사도 + pose 유사도 동시 고려
    \end{itemize}
    \item 결과: 해당 micro skill 내에서 test 객체에 가장 적합한 단일 시연 궤적 획득
  \end{itemize}
\end{frame}

%------------------------------------------------
\begin{frame}{검색 기반 정렬: Trajectory Transfer}
  \small
  \begin{itemize}
    \item 입력: test 장면의 타깃 객체 포인트 클라우드, 선택된 시연의 초기 EE pose
    \item 목표: 시연에서의 EE–객체 상대 pose를 test 장면에 보존하는 EE pose 계산
    \item 절차
    \begin{itemize}
      \item test 객체와 시연 객체 사이의 상대 pose $T_\delta$ 추정
      \item 식 $T^{\text{Test}}_{WE} = T_\delta T^{\text{Demo}}_{WE}$ 로 test EE pose 계산
      \item $T_\delta$: 회귀 기반 초기 추정 후 Generalized ICP로 정제
      \item motion planner로 해당 pose까지 경로 생성 및 실행
    \end{itemize}
    \item 특징: 정렬 단계에서 기하학적 추론과 경로 계획을 명시적으로 활용
  \end{itemize}
\end{frame}

%------------------------------------------------
\begin{frame}{검색 기반 상호작용: Open-loop Replay}
  \small
  \begin{itemize}
    \item 입력: 선택된 시연 궤적의 EE frame 속도 시퀀스
    \item 실행 방식
    \begin{itemize}
      \item 시연에서 기록된 EE 속도를 EE frame 기준으로 그대로 재생
      \item 그리퍼 개폐 상태도 시연과 동일하게 재현
    \end{itemize}
    \item 장점
    \begin{itemize}
      \item 복잡한 조작 dynamics를 학습할 필요 없이 시연 궤적 구조 보존
      \item 소량 시연에서도 안정적 상호작용 실행 가능
    \end{itemize}
    \item 한계
    \begin{itemize}
      \item 개방루프 – 실행 중 피드백 기반 보정 불가
      \item 초기 정렬 오차, grasp pose 차이, 환경 변화에 취약
    \end{itemize}
  \end{itemize}
\end{frame}

%------------------------------------------------
\begin{frame}{데이터셋 크기/다양성 실험: Micro skill 및 객체}
  \small
  \begin{itemize}
    \item 실험 1(데이터 크기): 4개 micro skill
    \begin{itemize}
      \item 관절 객체 조작, 변형체 조작, 퍼올리기(scooping), 삽입 등 포함
      \item 각 micro skill당 seen 3개, unseen 2개 작업 – 총 seen 12, unseen 8 작업
    \end{itemize}
    \item 실험 2(데이터 다양성): 10개 micro skill
    \begin{itemize}
      \item 실험 1의 4개 + 추가 6개 micro skill
      \item 각 다양성 설정에서 micro skill당 unseen 작업 2개 유지
    \end{itemize}
    \item 객체: 다양한 가방, 토스터, 상자, 팬 등 – 포인트 클라우드 기반 조작
  \end{itemize}
\end{frame}

%------------------------------------------------
\begin{frame}{데이터셋 크기/다양성 실험 설정}
  \small
  \begin{itemize}
    \item 실험 1: 작업당 시연 수 스케일링
    \begin{itemize}
      \item 4 micro skill, 총 20 작업(12 seen, 8 unseen)
      \item 작업당 시연 수 1 $\rightarrow$ 50까지 증가
      \item 50 시연/작업은 복잡 조작 학습에 충분한 수준으로 설정
    \end{itemize}
    \item 실험 2: 고정 시연 예산 하에서 작업 수 스케일링
    \begin{itemize}
      \item 총 시연 수 150 고정
      \item 10 작업(각 15 시연) $\rightarrow$ 30 작업(각 5 시연) $\rightarrow$ 50 작업(각 3 시연)
      \item 각 설정에서 micro skill당 unseen 작업 2개 포함(총 20 unseen)
    \end{itemize}
    \item 평가 프로토콜
    \begin{itemize}
      \item 각 task당 3회 rollout, micro skill 전체 평균
      \item 객체 위치: 80cm $\times$ 45cm 작업 공간 내 무작위
      \item yaw 방향 $\pm 180^\circ$ 회전 랜덤
    \end{itemize}
  \end{itemize}
\end{frame}

%------------------------------------------------
\begin{frame}{궤적 분해 및 정책 설계 그림}
  \small
  \begin{itemize}
    \item 좌측: monolithic vs decomposed 궤적 구조 비교
    \item 우측: 시스템 아키텍처 – 포인트 클라우드 + 언어 입력, BC와 retrieval 경로 구분
    \item 색상
    \begin{itemize}
      \item 보라: 멀티태스크 정책
      \item 파랑: retrieval 기반 구성요소
      \item 분홍: BC 기반 구성요소
    \end{itemize}
  \end{itemize}
  \vfill
  \begin{center}
    \includegraphics[width=0.6\textwidth, keepaspectratio]{extracted/figure_2.png}
    
    \scriptsize 궤적 분해와 정렬/상호작용 정책 조합 개요
  \end{center}
\end{frame}

%------------------------------------------------
\begin{frame}{데이터셋 크기/다양성에 따른 성능 개요}
  \small
  \begin{itemize}
    \item MT3(Ret-Ret): 모든 데이터 regime에서 가장 높은 성공률
    \item few-shot 영역(작업당 3 시연)에서 MT3가 50 시연을 사용하는 다른 모든 방법보다 우수
    \item 분해 기반 방법 전체(Ret-Ret, Ret-BC, BC-Ret, BC-BC)가 monolithic MT-ACT+보다 일관되게 높은 성능
    \item 정렬/상호작용 단계 모두에서 retrieval 사용 시 평균 성공률 상승
    \item monolithic BC는 데이터가 많고 작업 다양성이 클수록 스케일링 이점 나타남
  \end{itemize}
\end{frame}

%------------------------------------------------
\begin{frame}{데이터셋 크기 증가 효과}
  \small
  \begin{itemize}
    \item 모든 방법에서 시연 수 증가에 따라 seen/unseen 작업 성공률 상승
    \item retrieval 기반
    \begin{itemize}
      \item 더 많은 시연 후보 중에서 test 인스턴스에 더 적합한 시연 선택 가능
      \item 검색 품질 향상으로 일반화 성능 증가
    \end{itemize}
    \item BC 기반
    \begin{itemize}
      \item 더 넓은 공간/형상 변이 커버 – 표현 학습 및 공간 일반화 개선
    \end{itemize}
    \item 관찰된 추세
    \begin{itemize}
      \item 분해 기반: 1–10 시연 구간에서 급격한 성능 향상, 50 근처에서 plateau
      \item monolithic BC: 10–50 시연 구간에서 가속적 향상 – 충분한 데이터에서 분해 기반을 추월할 잠재력
    \end{itemize}
  \end{itemize}
\end{frame}

%------------------------------------------------
\begin{frame}{데이터셋 다양성 증가 효과}
  \small
  \begin{itemize}
    \item 고정 시연 예산(150) 하에서 작업 수 증가 → 작업당 시연 수 감소
    \item retrieval 기반
    \begin{itemize}
      \item unseen 작업: 더 많은 객체 인스턴스 덕분에 test 객체와 더 유사한 시연 검색 가능 → 성능 향상
      \item seen 작업: pose–geometry trade-off로 인해 성능 저하
      \begin{itemize}
        \item pose 유사도 vs 형상 유사도 사이의 근본적 균형 문제
      \end{itemize}
    \end{itemize}
    \item monolithic BC(MT-ACT+)
    \begin{itemize}
      \item 작업 다양성 증가가 seen/unseen 모두에 유리
      \item 10 micro skill, 50 작업, 150 시연 설정이 4 micro skill, 12 작업, 600 시연과 유사 성능 – 다양성을 데이터 효율성으로 전환
    \end{itemize}
    \item 분해 기반 BC-BC
    \begin{itemize}
      \item MT-ACT+와 동일 BC 구현 사용에도 다양성 이득 제한적
      \item 정렬/상호작용을 분리 학습함으로써 전체 궤적에서의 구조 공유가 약화된 영향
    \end{itemize}
  \end{itemize}
\end{frame}

%------------------------------------------------
\begin{frame}{데이터셋 크기/다양성 분석 그림}
  \small
  \begin{itemize}
    \item 상단 패널: 각 방법별 seen/unseen 성공률 vs 시연 수/작업 수
    \item 중간 패널: 분해 기반 평균 vs monolithic BC 비교 – 분해 prior의 전반적 이점
    \item 하단 패널: 정렬/상호작용 단계에서 BC vs retrieval 선택에 따른 평균 성능 차이
  \end{itemize}
  \vfill
  \begin{center}
    \includegraphics[width=0.6\textwidth, keepaspectratio]{extracted/figure_4.png}
    
    \scriptsize 데이터 크기와 다양성에 따른 성능 및 분해/검색 효과 분석
  \end{center}
\end{frame}

%------------------------------------------------
\begin{frame}{1000개 작업 학습: 설정 개요}
  \small
  \begin{itemize}
    \item 목표: MT3가 단일 시연만으로 매우 많은 실제 조작 작업을 학습 가능한지 검증
    \item 데이터
    \begin{itemize}
      \item 1000개 seen 작업 – 각 작업당 1개 시연
      \item 402개 객체, 31 macro skill, 534 micro skill 포함
      \item 추가 100개 unseen 작업 – 동일 macro skill 분포
      \item 시연 수집 시간: 단일 로봇에서 총 17시간
    \end{itemize}
    \item 평가
    \begin{itemize}
      \item seen/unseen 각 작업당 2회 rollout – 총 2200 rollout
      \item 5–20개 distractor, 다양한 조명, 위치/자세 랜덤(최대 $\pm 45^\circ$ yaw)
    \end{itemize}
    \item 목적: MT3의 강점/약점, open-loop replay의 근본 한계, 카테고리 수준 일반화 실패 모드 파악
  \end{itemize}
\end{frame}

%------------------------------------------------
\begin{frame}{1000개 작업 데이터셋 시각화}
  \small
  \begin{itemize}
    \item 상단: 시간 흐름에 따라 다양한 일상 조작 작업 프레임 나열 – 1000개 작업 스펙트럼 표현
    \item 하단: 사용된 객체 일부와 평가된 skill 유형 예시
    \item 포함 작업 예
    \begin{itemize}
      \item 문/서랍 열기, 물 붓기, 쌓기, 닦기, 삽입, 도구 사용 등
    \end{itemize}
  \end{itemize}
  \vfill
  \begin{center}
    \includegraphics[width=0.6\textwidth, keepaspectratio]{extracted/figure_1.png}
    
    \scriptsize 하루 만에 학습된 1000개 작업과 객체/스킬 분포 개략
  \end{center}
\end{frame}

%------------------------------------------------
\begin{frame}{1000개 작업: 예시 rollout 및 장면 다양성}
  \small
  \begin{itemize}
    \item 예시 rollout
    \begin{itemize}
      \item 팬에서 달걀 퍼올리기, 창문 닦기, 물 붓기 등 다양한 상호작용 시퀀스
    \end{itemize}
    \item 장면 다양성
    \begin{itemize}
      \item cluttered scene, 다양한 배경/조명, 여러 distractor 객체
      \item 타깃 객체 위치/자세 랜덤 – 실제 환경 유사 조건
    \end{itemize}
    \item 목적: MT3가 단순한 실험실 환경을 넘어 실제적 변이를 견딜 수 있는지 검증
  \end{itemize}
  \vfill
  \begin{center}
    \includegraphics[width=0.6\textwidth, keepaspectratio]{extracted/figure_5.png}
    
    \scriptsize 1000개 작업 평가에서의 rollout 예시와 장면 구성
  \end{center}
\end{frame}

%------------------------------------------------
\begin{frame}{1000개 작업: 전체 성능 및 macro skill별 결과}
  \small
  \begin{itemize}
    \item 전체 평균 성능
    \begin{itemize}
      \item seen 작업: 78.25\% 성공률
      \item unseen 작업: 68\% 성공률
    \end{itemize}
    \item macro skill별 패턴
    \begin{itemize}
      \item 접근 각도/접촉 위치 오차 허용도가 큰 작업에서 80\% 이상 높은 성공률
      \item 예: 닦기(wiping), 젓기(stirring), 놓기(placing), grasping 등
    \end{itemize}
    \item macro skill별 seen/unseen 작업 수는 불균형 – 일부 skill에서 분산 큼
  \end{itemize}
\end{frame}

%------------------------------------------------
\begin{frame}{1000개 작업: macro skill 성능 및 실패 요인 분석 그림}
  \small
  \begin{itemize}
    \item 상단 레이더 차트: macro skill별 seen/unseen 성공률 분포
    \item 하단 도넛 차트: 실패 원인 비율
    \begin{itemize}
      \item retrieval 실패, segmentation 실패, pose estimation 실패, kinematics/실행 실패 등
    \end{itemize}
  \end{itemize}
  \vfill
  \begin{center}
    \includegraphics[width=0.6\textwidth, keepaspectratio]{extracted/figure_6.png}
    
    \scriptsize 1000개 작업에서 macro skill별 성능과 주요 실패 모드 분해
  \end{center}
\end{frame}

%------------------------------------------------
\begin{frame}{공간 일반화 특성}
  \small
  \begin{itemize}
    \item 높은 허용도 작업에서 강점
    \begin{itemize}
      \item 닦기, 젓기, 놓기, grasping 등 – 접근 각도/접촉 위치 약간의 오차 허용
      \item open-loop replay가 작은 정렬 오차에도 성공 유지
    \end{itemize}
    \item 작은 비대칭 특징이 중요한 작업에서 한계
    \begin{itemize}
      \item 주전자 손잡이/주둥이, 카드 슬롯, 열쇠 구멍 등
      \item 포즈 추정이 전역 형상에 치우쳐 작은 특징을 무시 → 180도 회전 오차 등 발생
    \end{itemize}
    \item 고정밀 정렬 요구 작업
    \begin{itemize}
      \item 플러그 삽입, 작은 고리 걸기 등 – mm 수준 정렬 필요
      \item open-loop 상호작용이 정렬 오차를 보정하지 못해 실패율 증가
    \end{itemize}
  \end{itemize}
\end{frame}

%------------------------------------------------
\begin{frame}{카테고리 수준 일반화 특성}
  \small
  \begin{itemize}
    \item 성공 사례
    \begin{itemize}
      \item 객체 인스턴스 간에도 상호작용 궤적 구조가 유사한 경우
      \item 예: 다양한 표면 닦기, 접시 위에 머그컵 놓기, 머그 grasping 등
    \end{itemize}
    \item 실패 사례
    \begin{itemize}
      \item 인스턴스 기하 차이가 궤적에 큰 변화를 요구하는 경우
      \item 예: 주전자에서 다른 모양 컵으로 붓기 – 주둥이와 컵 가장자리 정렬 방식 변화
      \item 카드 리더기 스와이프 – 슬롯 위치가 인스턴스마다 달라 정밀 정렬 실패
    \end{itemize}
    \item retrieval의 구조적 한계
    \begin{itemize}
      \item 시연 간 ``보간'' 불가 – 항상 하나의 시연만 선택
      \item 필요한 궤적이 두 시연 사이에 위치할 때 적절한 중간 궤적 생성 불가
    \end{itemize}
    \item 변형체 조작
    \begin{itemize}
      \item 시각적 유사성만으로는 동역학 차이를 포착 불가
      \item 예: 책을 가방에 넣기 – 가방 플랩 강성/탄성 차이로 인해 시연 궤적 그대로는 실패 빈번
    \end{itemize}
  \end{itemize}
\end{frame}

%------------------------------------------------
\begin{frame}{Open-loop 상호작용의 근본 한계}
  \small
  \begin{itemize}
    \item 실행 중 온라인 보정/반응 제어 부재
    \begin{itemize}
      \item 궤적 시작 후 오류 감지 및 경로 수정 불가
      \item 실패 후 재시도는 가능하나 근본적 해결 아님
    \end{itemize}
    \item 변형체 조작
    \begin{itemize}
      \item 천 접기, 옷 정리 등 – 재료 반응에 따라 연속적 조정 필요
      \item 고정 궤적 재생으로는 다양한 동역학에 대응 불가
    \end{itemize}
    \item 반응적 조작 요구 작업
    \begin{itemize}
      \item 접촉 기반 재배치, 다단계 밀기, 장애물 회피 등
      \item 지속적 피드백과 적응이 필수 – open-loop replay로는 원천적으로 불가능
    \end{itemize}
  \end{itemize}
\end{frame}

%------------------------------------------------
\begin{frame}{실패 모드 정량 분석}
  \small
  \begin{itemize}
    \item 분석 대상: seen 작업 실패 사례
    \item 평가 항목
    \begin{itemize}
      \item segmentation 정확성
      \item retrieval 정확성
      \item pose estimation 성공 여부
      \item motion execution/kinematics 문제
    \end{itemize}
    \item 주요 실패 비율
    \begin{itemize}
      \item retrieval 실패: 22.3\% – 부분 가림, 작은 특징 차이 등에서 전역 매칭 한계
      \item segmentation 실패: 19.5\% – 투명체, cluttered scene 등
      \item pose estimation 실패: 23.9\% – 비대칭 객체에서 큰 pose 차이 시 부분 포인트 클라우드 불일치
      \item grasped object 관련 실패: 20.2\% – 시연과 배치 시 grasp pose 불일치가 궤적 전체에 전파
    \end{itemize}
  \end{itemize}
\end{frame}

%------------------------------------------------
\begin{frame}{Retrieval 기반 방법의 해석 가능성}
  \small
  \begin{itemize}
    \item 정렬 단계
    \begin{itemize}
      \item 등록된 포인트 클라우드 중첩 시각화로 pose estimation 결과 직접 확인 가능
      \item misregistration 사전 감지 후 실행 중단/재시도 용이
    \end{itemize}
    \item 상호작용 단계
    \begin{itemize}
      \item 실행 궤적이 항상 시연 데이터 버퍼에서 선택된 궤적에 한정
      \item 행동이 ``보여준 대로''에서 벗어나지 않음 – 예측 가능성 높음
    \end{itemize}
    \item BC 대비 장점
    \begin{itemize}
      \item BC는 결정이 네트워크 가중치 내부에 은닉 – 해석 어려움
      \item retrieval는 어떤 시연이 선택되었는지, 어떤 pose가 사용되는지 명시적
    \end{itemize}
    \item 새로운 작업 추가
    \begin{itemize}
      \item 단순히 시연을 데이터셋에 추가 – 재학습/파인튜닝 불필요
      \item 빈번한 신규 작업 도입 환경에 적합
    \end{itemize}
  \end{itemize}
\end{frame}

%------------------------------------------------
\begin{frame}{Retrieval 기반 구성요소 시각화}
  \small
  \begin{itemize}
    \item 좌측: pose estimation 결과 – 시연/테스트 포인트 클라우드 정합 상태 시각화
    \item 우측: 상호작용에서 BC vs retrieval 사용 시 궤적 안정성 비교
    \begin{itemize}
      \item retrieval: 시연 궤적에 밀착된 안정적 경로
      \item BC: 데이터 부족 시 궤적 변동성 증가
    \end{itemize}
  \end{itemize}
  \vfill
  \begin{center}
    \includegraphics[width=0.6\textwidth, keepaspectratio]{extracted/figure_7.png}
    
    \scriptsize pose estimation 해석 가능성과 상호작용 안정성 비교 예시
  \end{center}
\end{frame}

%------------------------------------------------
\begin{frame}{BC의 데이터 요구와 스케일링 특성}
  \small
  \begin{itemize}
    \item low-data regime에서 BC의 어려움
    \begin{itemize}
      \item 포인트 클라우드 $\rightarrow$ 궤적 재배치/제어 명령 매핑 학습
      \item 언어 조건부 멀티태스크 구분 학습
      \item 큰 workspace 내 위치/자세 변이, 카테고리 내 형상 변이 모두 일반화 필요
    \end{itemize}
    \item 데이터 부족 시
    \begin{itemize}
      \item 의미 있는 패턴 대신 시연 memorization 경향
      \item 새로운 pose/객체 인스턴스에 대한 일반화 실패
    \end{itemize}
    \item high-data regime에서의 잠재력
    \begin{itemize}
      \item 전체 궤적에서 alignment/interaction 패턴을 함께 학습
      \item 다양한 작업 간 공유 구조를 활용해 표현 학습
      \item 충분한 시연과 다양성 제공 시 MT3보다 우수해질 가능성
    \end{itemize}
  \end{itemize}
\end{frame}

%------------------------------------------------
\begin{frame}{연구 한계 및 적용 범위}
  \small
  \begin{itemize}
    \item 단일 상호작용, 단일 팔 작업에 초점
    \item grasped object pose는 시연/배치 간 동일하다고 가정
    \begin{itemize}
      \item 기존 방법을 통해 추가 시연 없이 grasp pose 차이 보정 가능하다는 전제
    \end{itemize}
    \item 환경 제약
    \begin{itemize}
      \item 시연 궤적을 방해하는 장애물/동적 객체는 고려하지 않음
      \item open-loop replay는 장애물 무시 → 충돌 가능
      \item BC는 out-of-distribution 상황 – 추가 시연 필요
    \end{itemize}
    \item 감각 모달리티
    \begin{itemize}
      \item 비전만 사용 – 변형체 동역학(강성, 마찰 등) 추론 불가
      \item 촉각/힘 센서 미사용
    \end{itemize}
    \item 다관절/다단계 작업
    \begin{itemize}
      \item bimanual, multi-step 작업은 고수준 planner/skill chaining으로 확장 가능하나 본 연구에서 직접 실험하지 않음
    \end{itemize}
  \end{itemize}
\end{frame}

%------------------------------------------------
\begin{frame}{결론 요약}
  \small
  \begin{itemize}
    \item 궤적 분해 + 검색 기반 일반화가 few-shot imitation learning에서 강력한 대안임
    \item MT3
    \begin{itemize}
      \item 작업당 1–3 시연에서도 높은 성공률 달성
      \item 1000개 일상 조작 작업을 24시간 이내 시연으로 학습
      \item seen 78.25\%, unseen 68\% 성공률로 실질적 유용성 입증
    \end{itemize}
    \item 분해 prior
    \begin{itemize}
      \item low-data regime에서 monolithic BC 대비 한 order 수준 데이터 효율성 향상
      \item 정렬/상호작용에 특화된 정책으로 구조적 inductive bias 제공
    \end{itemize}
    \item 접근 선택 기준
    \begin{itemize}
      \item 데이터/시연 비용이 제한된 응용: MT3와 같은 분해+retrieval 방식 적합
      \item 대규모 시연 수집 가능, 높은 다양성 확보 가능: monolithic BC가 장기적으로 유리
    \end{itemize}
  \end{itemize}
\end{frame}

%------------------------------------------------
\begin{frame}{향후 연구 방향 1: MT3 개선}
  \small
  \begin{itemize}
    \item 검색 모듈 고도화
    \begin{itemize}
      \item 부분 가림/작은 특징에 강인한 기하 표현 학습
      \item pose–geometry trade-off를 명시적으로 다루는 다목적 최적화 기반 retrieval
      \item 다수 시연을 활용한 궤적 ``보간'' 또는 혼합 전략
    \end{itemize}
    \item open-loop 한계 보완
    \begin{itemize}
      \item 시연 궤적을 초기화로 사용하고, 실행 중 폐루프 제어로 미세 조정
      \item 실패 감지 및 온라인 replanning 메커니즘 통합
    \end{itemize}
    \item grasp pose 불일치 대응
    \begin{itemize}
      \item self-supervised grasp adaptation 기법과 MT3 통합
      \item grasp-invariant 표현을 활용한 trajectory transfer
    \end{itemize}
  \end{itemize}
\end{frame}

%------------------------------------------------
\begin{frame}{향후 연구 방향 2: 범용성 확장}
  \small
  \begin{itemize}
    \item 변형체 및 동역학 풍부한 작업
    \begin{itemize}
      \item 촉각/힘 센서 통합으로 재료 특성 추론
      \item 동역학-aware trajectory transfer 또는 hybrid RL-IL 구조 탐구
    \end{itemize}
    \item 복잡한 환경 및 장애물 존재 시나리오
    \begin{itemize}
      \item 시연 궤적을 제약으로 사용하는 모션 플래닝 결합
      \item 환경 변화에 강인한 local replanning 전략
    \end{itemize}
    \item bimanual 및 multi-step 작업
    \begin{itemize}
      \item dual-arm alignment/interaction 분해 구조 설계
      \item LLM 기반 high-level planner와 MT3 primitive의 조합
    \end{itemize}
  \end{itemize}
\end{frame}

%------------------------------------------------
\begin{frame}{향후 연구 방향 3: BC와의 하이브리드}
  \small
  \begin{itemize}
    \item retrieval + BC 결합 전략
    \begin{itemize}
      \item retrieval로 선택한 시연을 context로 사용하는 conditional BC
      \item retrieval로 초기 궤적 제공, BC가 fine-tuning 역할 수행
    \end{itemize}
    \item 데이터 스케일에 따른 적응형 프레임워크
    \begin{itemize}
      \item low-data: MT3 중심, high-data: monolithic BC 비중 증가
      \item 데이터/작업 분포에 따라 분해 수준과 학습 방식 자동 선택
    \end{itemize}
    \item 대규모 웹/시뮬레이션 데이터 활용
    \begin{itemize}
      \item 웹 비디오/시뮬레이션에서 사전 학습한 BC 정책 위에 MT3-style retrieval layer 추가
      \item VLA 모델과의 결합으로 언어/시각/행동 통합 일반화 강화
    \end{itemize}
  \end{itemize}
\end{frame}

%------------------------------------------------
\begin{frame}{요약 및 시사점}
  \small
  \begin{itemize}
    \item 구조적 prior(정렬/상호작용 분해)와 시연 검색만으로도 대규모 작업 학습 가능함
    \item 복잡한 거대 신경망 정책 없이도 1000개 작업 수준의 스케일 달성 가능함
    \item 로봇 조작 연구에서 ``데이터 효율성 vs 표현력''의 새로운 Pareto front 제시
    \item 실제 응용에서
    \begin{itemize}
      \item 공장/서비스 로봇: 작업 변경이 잦고 시연 비용이 큰 환경에 MT3 유리
      \item 장기적으로는 MT3와 대규모 BC/VLA 모델의 통합이 범용 로봇 학습의 핵심 방향으로 기대됨
    \end{itemize}
  \end{itemize}
\end{frame}

%------------------------------------------------
\end{document}